% Source: https://artofproblemsolving.com/community/c3h1403898_math_kangaroo_2006_usa_78_pb_30
% Stable identity: aops-topic-1403898-q30
Each of (A)-(D) has a counterexample:
\begin{itemize}
\item (A): $x=y=10$, $z=0.1$ gives $xy=100$.
\item (B): $x=20$, $y=z=0.05$ gives $xy=1$.
\item (C): $x=10$, $y=7.5$, $z=2.6$ gives $xy=75$.
\item (D): $x=\dfrac{125}{7}$, $y=\dfrac{7}{5}$, $z=20.1-\dfrac{125}{7}-\dfrac{7}{5}=\dfrac{1407-1250-98}{70}=\dfrac{59}{70}$. Then $xy=25$, and $x\ge y\ge z>0$ since $\dfrac{125}{7}>\dfrac{7}{5}=\dfrac{98}{70}>\dfrac{59}{70}$.
\end{itemize}
(For (D), set $xy=25$, $x=\dfrac{25}{y}$; then $z=20.1-\dfrac{25}{y}-y$, and $0<z\le y$ holds for $y$ roughly between $1.33$ and $1.45$.) So none of (A)-(D) is always true, and the answer is $\boxed{\text{(E)}}$.
